\scp{\tsc{Seram-Tanimbar-Bomberai} classification}{11-06}
In addition to the changes listed below,
nearly all STB languages have *p > \it{f}.
The only exceptions are Kur with *p = \it{p} in all word positions
and Yamdena with word-final *p = \it{p}.

\dirtree{%%The percent after { is needed
.1 {\tsc{Seram-Tanimbar-Bomberai} *z/*d > **d; *ə/*a/*-aw/*-ay > **a /σ\und{σ}{\#}}.
		.2 {\tsc{East Seram} *ŋ > \it{n-, -k-}; *ma-p > \it{b}; *k > {\0}, \it{(ʔ)} [\it{k-} some prefixes]; \newline\hp{**}*d/*z/*l/*j > \it{l}; *u/*i/*uy > **u /σ\und{σ}{\#}; **u > \it{i} /{\gap}C{\#}; *{\0} > \it{y} /{\#}{\gap}a [LS]}.
				.3 {Bobot [bty] *b > \it{v}; *R > \it{w}, {\0}; **u > {\0}, u̥ /{\gap}{\#}}.
				.3 {Masiwang (a.k.a. Bonfia) [bnf] *s/*R > \it{s}; **u > \it{i} /{\gap}{\#}}.
		.2 {\tsc{Eastern Islands} *j > **s, **l; *p > **f; *map > **maf (*q = *q; *ŋ = *ŋ)}.
				.3 {\tsc{Nuclear Eastern Islands} *b > \it{w}; *b > {\0}, \it{w} /{\gap}u}.
						.4 {\tsc{Geser-Gorom}}.
								.5 {Bati [bvt] **d/*l > \it{l} {\#}{\gap}; **d/*l > \it{l,} {\0} /V{\gap}V; *R > \it{l} (*k = \it{k})}.
								.5 {\tsc{Nuclear Geser-Gorom} [ges] **d/*R > \it{r}}.
										.6 {Geser **f > \it{h}; *k > \it{ʔ}, {\0}}.
										.6 {Gorom}.
						.4 {\tsc{Watubela} *s/**s > \it{s, h}; **d/*r > **r}.
								.5 {Watubela [wah] *q/*k > \it{k}; **r/*l > \it{l}}.
								.5 {Kesui [{-}{-}{-}] *q/*k > {\0}, \it{h}; (**r = \it{r})}.
				.3 {Kowiai [kwh] *b/**f > \it{ɸ}; **d/*R/*l > \it{r}; *k > \it{ʔ}, {\0}; *ŋ > \it{g, ŋg}}.
		.2 {\tsc{Teor-Kur} *j > {\0}; *b > \it{f} ; **d/*R > \it{r}; *s > \it{s,} **h; **a > \it{i} /σ\und{σ}{\#} \tsc{3sg poss}}.
				.3 {Teor [tev] *ŋ > \it{g}; **h > {\0}; *k > \it{k, h}; **a > ‹a›, ‹e› /iC{\gap}(C){\#}}.
				.3 {Kur [kuv] (*p = \it{p}; *k = \it{k})}.
		.2 {\tsc{Tanimbar-Bomberai} *j > **j, \it{r}; *R > \it{r} (\srf{11-01-01})}.
				%.3 {Yamdena [jmd] **d- = \it{d /{\#}{\gap};} **-d- > -\it{r-, -d-} /V{\gap}V; *u > \it{i} /σ\und{σ}{\#}; *ә > \it{e} /\und{σ}σ{\#} (*k = \it{k}; *ŋ = \it{ŋ}; *-p = -\it{p} /{\gap}{\#})}.
				%.3 {\tsc{Nuclear Tanimbar-Bomberai} **d > **l /{\gap}Vr; **d/**j > **r /else.; *{\0} > \it{y} /{\#}{\gap}V}.
						%.4 {\tsc{Kei-Fordata} *b > \it{v}; *ŋ > \it{n}; *s > **h; (**l = \it{l})}.
								%.5 {\tsc{Fordata} [frd] **h > \it{h, \0}}.
										%.6 {Fordata}.
										%.6 {Sera}.
								%.5 {Kei [kei]}.
										%.6 {Kei Kecil **d/**r > \it{r}}.
										%.6 {Kei Besar **r > \it{x}; **d > \it{r}}.
						%.4 {\tsc{Bomberai Tip} **d/**l/**r > \it{n}; *u > \it{i} /σ\und{σ}{\#}}.
								%.5 {Uruangnirin [urn] *y > \it{l}; *b > \it{p}; *ŋ > \it{g, ŋg, ŋ}}.
								%.5 {\tsc{Onin-Sekar} **n > \it{r} /{\gap}Vr; *ŋ > \it{g}}.
										%.6 {Onin [oni] *b > \it{p}}.
										%.6 {Sekar [skz]}.
				%.3 {\tsc{East Bomberai} *l > \it{r}; *b > \it{p}; *ŋ > **g; *V > {\0} /{\gap}{\#}; *{\0} > **ə /C{\gap}{\#}}.
						%.4 {\tsc{Nuclear East Bomberai} lexically specific vowel changes}.
								%.5 {Arguni [agf] **ə > \it{e}}.
								%.5 {\tsc{Bedoanas-Erokwanas} **g > \it{ʔ}; **ə > \it{a}}.
										%.6 {Bedoanas (a.k.a. Yarik) [bed] *p > {\0}}.
										%.6 {Erokwanas (a.k.a. Kambran, Damberang-Goras) [erw]}.
						%.4 {\tsc{Irarutu-Kuri} *k > {\0}; **d/*l/**r > \it{r}; *s > \it{s, {\0}}}.
								%.5 {Irarutu [irh]}.
								%.5 {Kuri (a.k.a. Nabi) [nbn]}.
}

\newpage
\sscp{\tsc{Tanimbar-Bomberai} internal structure}{11-06-01}
\dirtree{%%The percent after { is needed
.1 {\tsc{Seram-Tanimbar-Bomberai} *z/*d > **d; *ə/*a/*-aw/*-ay > **a /σ\und{σ}{\#}}.
		.2 {\tsc{Tanimbar-Bomberai} *j > **j, \it{r}; *R > \it{r}}.
				.3 {Yamdena [jmd] **d- = \it{d /{\#}{\gap};} **-d- > -\it{r-, -d-} /V{\gap}V; *u > \it{i} /σ\und{σ}{\#}; *ә > \it{e} /\und{σ}σ{\#} (*k = \it{k}; *ŋ = \it{ŋ}; *-p = -\it{p} /{\gap}{\#})}.
				.3 {\tsc{Nuclear Tanimbar-Bomberai} **d > **l /{\gap}Vr; **d/**j > **r /else.; \newline\hp{**}*{\0} > \it{y} /{\#}{\gap}V}.
						.4 {\tsc{Kei-Fordata} *b > \it{v}; *ŋ > \it{n}; *s > **h; (**l = \it{l})}.
								.5 {\tsc{Fordata} [frd] **h > \it{h, \0}}.
										.6 {Fordata}.
										.6 {Sera}.
								.5 {Kei [kei]}.
										.6 {Kei Kecil **d/**r > \it{r}}.
										.6 {Kei Besar **r > \it{x}; **d > \it{r}}.
						.4 {\tsc{Bomberai Tip} **d/**l/**r > \it{n}; *u > \it{i} /σ\und{σ}{\#}}.
								.5 {Uruangnirin [urn] *y > \it{l}; *b > \it{p}; *ŋ > \it{g, ŋg, ŋ}}.
								.5 {\tsc{Onin-Sekar} **n > \it{r} /{\gap}Vr; *ŋ > \it{g}}.
										.6 {Onin [oni] *b > \it{p}}.
										.6 {Sekar [skz]}.
				.3 {\tsc{East Bomberai} *l > \it{r}; *b > \it{p}; *ŋ > **g; *V > {\0} /{\gap}{\#}; *{\0} > **ə /C{\gap}{\#}}.
						.4 {\tsc{Nuclear East Bomberai} (lexically specific vowel changes)}.
								.5 {Arguni [agf] **ə > \it{e}}.
								.5 {\tsc{Bedoanas-Erokwanas} **g > \it{ʔ}; **ə > \it{a}}.
										.6 {Bedoanas (Yarik) [bed] *p > {\0}}.
										.6 {Erokwanas (Kambran, Damberang-Goras [erw]}.
						.4 {\tsc{Irarutu-Kuri} *k > {\0}; **d/*l/**r > \it{r}; *s > \it{s, {\0}}}.
								.5 {Irarutu [irh]}.
								.5 {Kuri (Nabi) [nbn]}.
}
